我叫马睿骁，目前为浙江理工大学物理学专业硕士研究生，预计于2027年6月毕业。硕士阶段，我主要围绕\textbf{低维功能材料、多铁耦合及相关量子物性}开展第一性原理计算和多尺度模拟研究。

随着科研工作的深入，我的关注重点逐渐从“如何完成一个计算”转向“如何真正解释一个物理问题”。相比单纯预测材料是否具有某种性质，我越来越关注原子结构、极化、轨道杂化、界面效应与材料功能之间的内在联系，也逐渐形成了对\textbf{铁电/反铁电材料、极化翻转、外延薄膜和异质界面}等问题的研究兴趣。

硕士阶段的工作使我积累了较扎实的第一性原理、科研编程、数值模拟和机器学习基础。博士阶段，我希望继续发挥已有的计算与程序开发优势，同时系统补充材料制备和实验表征能力，逐步形成能够连接\textbf{理论预测、数值模拟与实验验证}的研究能力。

\section{一、专业基础与研究特长}

硕士期间，我接受了较系统的凝聚态物理和计算材料研究训练。目前较为熟悉 \textbf{VASP、Wannier90/WannierTools、TB2J、Spirit} 等计算工具，能够开展结构优化、电子结构、自旋轨道耦合、铁电极化、极化翻转路径、界面结构、磁交换作用、磁各向异性、Berry曲率及相关拓扑性质等计算，并根据不同问题进一步建立有效模型和开展数值模拟。

在长期科研过程中，我逐渐形成了一套相对完整的研究思路：

\textbf{从晶体结构出发，分析电子结构和轨道行为，再进一步追踪微观相互作用，最终理解材料宏观功能性质的来源。}

相比软件本身，我更重视计算结果是否可靠以及不同证据之间能否相互支持。因此，对于 U 值、SOC、k点密度、Wannier拟合窗口、有限尺寸效应等可能影响结论的因素，我通常会主动进行参数测试和对照计算。当结果出现异常时，我也更倾向于重新检查结构、计算设置和模型假设，而不是直接选择符合预期的结果。

\subsection*{编程与科研代码开发}

\textbf{编程是我目前比较突出的研究能力之一。}

我能够使用 \textbf{Python、Shell/Linux} 独立编写科研程序和自动化脚本，用于高性能计算任务提交、批量模型生成、参数扫描、结果提取、数据处理、统计分析和可视化。对我而言，编程并不仅仅是提高计算效率，更重要的是它使我能够根据新的物理问题自行建立分析方法，而不完全受限于现有软件已经提供的功能。

一个较为典型的例子来自 TcIrGe₂Se₆ 的研究。为了理解其中的长程交换机制，仅依靠能带和投影态密度难以定量判断不同轨道通道的贡献。为此，我基于 Wannier 紧束缚模型自行编写 Python 程序，解析 \texttt{wannier90\_hr.dat}、Wannier中心和晶体结构信息，并进一步提取和统计不同原子、不同轨道之间的跃迁矩阵元。

通过这一方法，我能够从定量角度分析特定交换路径中的轨道杂化，将原本主要依靠图像和经验判断的问题转化为可以计算和比较的数据。

这类经历使我逐渐形成了一个科研习惯：\textbf{如果现有方法不能直接回答问题，就进一步理解方法本身，再尝试通过程序把问题转化为可以计算和验证的形式。}

这也是我希望博士阶段继续发挥的优势。

\section{二、科研经历与创新实践}

硕士期间，我主要围绕二维多铁材料及低维功能材料开展研究。

在二维多铁材料 \textbf{TcIrGe₂Se₆} 的研究中，我首先分析材料的结构稳定性和铁电极化，在此基础上进一步研究不同极化状态下的电子结构、磁性、谷自由度和Berry曲率等性质。

为了理解其中微观相互作用的来源，我进一步构建 Wannier 紧束缚模型，对 \textbf{Tc–Se–Ir–Se–Tc 长程超超交换通道}进行分析，将局域结构变化、轨道杂化和微观交换作用联系起来。

这一研究使我逐渐建立了从

\textbf{结构变化 → 极化状态 → 电子结构 → 轨道杂化 → 微观相互作用 → 功能性质}

出发理解材料的研究思路。

相关工作 \textbf{“Ferroelectric control of magnetism, valley polarization, and skyrmions in monolayer TcIrGe₂Se₆”} 已被 \textit{Physical Review B} 接收。

除上述工作外，我还开展了成分替位调控二维多铁材料性质的研究，并将第一性原理获得的参数进一步用于有限温度模型分析。目前，我正在推进二维铁电/磁性异质结构相关研究，主要关注\textbf{极化翻转、界面构型、层间耦合和应变对材料性质的影响}。

这些研究经历也逐渐改变了我的研究兴趣。相比不断寻找具有某一种特殊性质的新材料，我现在更希望回答：

\textbf{材料性质为什么发生变化？这种变化能否通过结构、极化、界面或外场进行主动控制？}

这是我希望博士阶段继续深入的问题。

\section{三、机器学习与新方法探索}

除了传统第一性原理和模型计算，我也在主动探索\textbf{机器学习与物理问题的结合}。

在交错磁拓扑研究中，我尝试将 \textbf{Sobol参数采样、拓扑不变量计算、机器学习预测和局部参数搜索}结合起来。

面对较大的模型参数空间，如果完全依赖人工逐点扫描，不仅计算量较大，也很难迅速判断真正值得进一步分析的区域。因此，我先利用数值方法生成参数数据，再使用机器学习模型辅助判断可能出现拓扑相变的区域，最后重新回到严格的物理计算中进行确认。

这一研究让我对机器学习形成了比较明确的认识：

\textbf{机器学习不应该替代物理，而应该帮助研究者更高效地发现物理。}

我更希望利用它处理传统方法难以高效完成的参数搜索、材料筛选和复杂数据分析问题，而最终的物理结论仍然需要建立在严格计算和机制分析之上。

未来，我尤其希望进一步学习\textbf{机器学习原子势}。如果能够以第一性原理数据训练原子势，就有机会把目前DFT能够处理的时间和空间尺度进一步拓展，用于研究更大体系中的极化翻转、畴壁运动、缺陷作用以及有限温度结构演化。

我认为这能够很好地连接我目前已有的\textbf{第一性原理、编程和机器学习能力}，并服务于博士阶段计划研究的铁电材料问题。

\section{四、研究兴趣与研究设想}

经过硕士阶段的研究，我目前最希望继续深入的是\textbf{铁电/反铁电材料、外延薄膜及异质界面中的结构与极化问题}。

首先，我非常关注\textbf{极化翻转过程本身}。

很多计算研究最终比较的是两个稳定极化态，但真实极化翻转往往涉及原子的连续移动、中间亚稳结构、局域结构变化甚至缺陷参与。我希望进一步理解不同翻转路径为什么具有不同能垒，以及应变、界面和缺陷如何影响极化翻转。

其次，我对\textbf{铁电相、反铁电相和其他亚稳结构之间的竞争}很感兴趣。

在真实材料中，不同结构之间的稳定性往往受到晶向、外延应变、电场、界面和缺陷等多种因素共同影响。这意味着通过材料设计和外部条件，有可能人为改变材料所处的能量状态和相变路径。我希望结合第一性原理、NEB及分子动力学等方法研究这一过程。

此外，我也希望进一步研究\textbf{外延薄膜和异质界面}。

相比理想自由材料，真实薄膜中的外延约束、界面结构、电荷重构和缺陷会带来更加复杂但也更加接近实际材料的问题。这样的体系既有基础物理价值，也更加容易与实验建立直接联系。

对于极化畴、畴壁、极化涡旋及其他复杂极化结构，我同样具有较强兴趣。这些结构涉及晶格、静电作用、弹性场和界面条件之间的竞争，也是连接原子尺度结构和宏观功能响应的重要问题。

\section{五、博士阶段学习与研究计划}

博士阶段前期，我希望进一步加强凝聚态物理、铁电物理和材料相变等方面的理论基础，并尽快熟悉所在课题组的材料体系和实验方法。

在计算方面，我会继续发挥已有的第一性原理和程序开发能力，并进一步学习分子动力学、机器学习原子势以及适用于铁电体系的多尺度模拟方法。

实验能力将是我博士阶段希望重点补充的部分。

我希望系统学习\textbf{外延薄膜生长、铁电性能测试、PFM、高分辨STEM及同步辐射结构表征}等实验方法，真正了解理论模型中的结构、极化和相变如何在真实样品中表现出来。

博士阶段中期，我希望能够围绕具体材料逐渐形成相对独立的研究课题，并尝试真正将理论计算与实验结合。

例如，可以从计算出发研究应变、界面或缺陷可能导致的结构变化，再通过实验进行观察；反过来，如果实验中出现理论模型无法直接解释的现象，也可以重新建立结构模型和计算方案寻找原因。

我希望最终形成的不是简单的“计算预测、实验验证”，而是：

\textbf{计算帮助实验提出问题，实验帮助理论发现新的问题。}

博士阶段后期，我希望逐步提高独立凝练科学问题、设计研究方案和推进完整课题的能力，在完成博士论文和相关研究成果的基础上，逐渐形成自己的研究特色。

\section{六、职业规划}

对于未来职业选择，我希望保持一定的开放性。

如果博士阶段能够取得较好的科研成果，并有机会进入合适的平台继续开展研究，我愿意在博士毕业后继续进行博士后训练，进一步积累研究成果、合作经验和独立科研能力，并争取进入高校或科研院所。

高校教师和科研人员仍然是我比较向往的发展方向，因为这样的职业能够让我持续接触新的问题并保持较长期的研究积累。

但我也认为，职业规划需要根据博士阶段的实际情况不断调整。

如果届时自己的研究成果、发展机会或个人条件更适合其他道路，我也会认真考虑高校教学岗位、科研院所或企业研发等选择。尤其在新材料、半导体、计算材料和人工智能辅助材料研发等方向，企业同样能够提供较好的研究和技术平台。

因此，相比现在就限定自己未来必须成为哪一种职业，我更希望在博士阶段不断提升自己的能力和科研竞争力，使自己在未来拥有更多选择。

如果科研成果、研究平台和发展机会能够支持我继续向前，我愿意继续走学术道路；如果实际情况更适合其他方向，我也会选择能够发挥自己专业能力并继续成长的岗位。

无论最终走向高校、科研机构还是企业，我希望自己始终能够保持对新问题的兴趣，并尽量从事具有一定研究性和创造性的工作。

\section{七、寄语——从理论预测走向真实材料}

虽然我的硕士阶段主要从事理论计算，但相比长期停留在纯理论计算中，我其实更加向往能够亲自参与实验。

第一性原理计算为我理解材料提供了非常重要的工具，也构成了我目前最主要的科研基础。但几年的计算研究也让我逐渐意识到理论模型自身的局限。

计算结果往往受到结构模型、计算方法、交换关联参数和数值设置等因素影响，而真实实验样品还会存在温度、缺陷、应变、界面和动力学过程等更加复杂的条件。

因此，我越来越想真正回答这样一个问题：

理论上预测出来的材料和物性，在实验中究竟能不能被做出来、被看到，并最终被测量到？

这是我非常向往的一种科研状态。

如果能够先从理论上提出一个材料或结构，通过计算判断其稳定性和可能具有的性质，再亲自参与材料制备与实验表征，最终在实验中观察到与理论预测相对应的现象，我认为这种过程会带给我非常直接的科研反馈。

对我而言，这种反馈不仅意味着“计算结果得到了验证”，更重要的是说明自己的研究真正和现实材料建立了联系。

理论预测可以帮助实验减少盲目试错，而实验结果又可以反过来检验甚至修正理论模型。我希望自己的研究能够接受真实实验的检验，而不仅仅停留在一个理想化的模型中。

如果实验结果和理论预测一致，我会非常享受这种从理论到现实的过程；如果两者并不一致，我同样愿意继续追问其中的原因。因为很多真正有意义的问题，往往就隐藏在理论模型与真实材料之间的差别中。

这也是我博士阶段尤其希望学习实验的重要原因。

我并不希望简单地从理论“转行”到实验，也不认为理论和实验是彼此替代的研究方式。相反，我更希望最终能够站在两者之间：既可以利用第一性原理和数值方法预测材料、理解微观机制，也能够理解材料如何制备、实验信号如何获得，以及理论中的理想体系与真实材料之间究竟存在怎样的差异。

我很喜欢这样一种研究过程：

提出一个想法， 用理论判断它是否可能实现， 再尝试真正把它做出来， 最后通过实验判断最初的想法是否成立。

这种从“预测”走向“实现”的过程，是我目前对科研最向往的状态之一。

同时，我也一直希望有机会进入研究条件更加丰富、学术交流更加充分的平台。对我而言，更高水平的平台不仅意味着更多实验设备和计算资源，更重要的是能够接触不同研究背景的人、更加前沿的问题以及更多自己以前没有使用过的方法。

我希望知道自己在这样的环境中还能够学会什么，也希望看到自己在科研道路上究竟能够走多远。

探索未知一直是一件让我感到有趣的事情。如果未来的科研环境、个人发展以及现实条件允许，我希望能够长期从事科研工作。相比已经知道答案的问题，我始终更容易被那些“为什么会这样”“是否还存在另外一种可能”“理论预测能否真正实现”的问题吸引。

因此，对我而言，博士阶段不仅是一次学历上的继续深造，也是一次重新拓展研究边界的机会。

我希望自己能够从目前以理论计算为主要基础的硕士生，逐渐成长为能够在理论预测、程序开发、数据分析、材料制备、实验表征和物理解释之间建立联系的研究者，并在不断探索未知的过程中，找到适合自己长期发展的方向。
